% !TeX program = platex
% UTF-8. Compile with pLaTeX (twice) and dvipdfmx.
% All drawing elements are editable TikZ vectors; no external images are required.
% Dimension interpretation used in this drawing:
%   projection normal to the front face = 3 cm
%   vertical height of the grip/ledge   = 5 cm
%   ledge underside to body lower edge  = 15 cm
%   long-axis length L                  = sufficiently long; not specified numerically
%   fillet radius of the tip upper edge  = 5 mm (edge_radius in the simulation model)
% The unspecified upper height, body depth, supports and hand are schematic.
\documentclass[a4paper,10pt,dvipdfmx]{jarticle}
\usepackage[landscape,margin=12mm]{geometry}
\usepackage[dvipdfmx]{graphicx}
\usepackage{tikz}
\usetikzlibrary{arrows.meta,calc,patterns}
\usepackage[dvipdfmx,hidelinks]{hyperref}
\usepackage{pxjahyper}
\usepackage[haranoaji]{pxchfon}
\renewcommand{\familydefault}{\sfdefault}
\renewcommand{\kanjifamilydefault}{\gtdefault}
\hypersetup{pdftitle={クリフの正面図と断面図},pdfauthor={d-hacks B3 omote}}
\pagestyle{empty}
\setlength{\parindent}{0pt}
\setlength{\parskip}{0pt}
\definecolor{Ink}{HTML}{24303B}
\definecolor{Slate}{HTML}{69737D}
\definecolor{Line}{HTML}{CBD1D6}
\definecolor{Body}{HTML}{F0F2F4}
\definecolor{Ledge}{HTML}{245B76}
\definecolor{LedgeFill}{HTML}{DCEAF0}
\definecolor{Hand}{HTML}{A76A3E}
\definecolor{HandFill}{HTML}{F6EDE6}
\definecolor{Contact}{HTML}{9E4D3F}
\tikzset{
  edge/.style={draw=Ink,line width=.65pt},
  dim/.style={<->,>=Latex,line width=.5pt,draw=Ink},
  ext/.style={draw=Slate,line width=.35pt},
  leader/.style={draw=Slate,line width=.45pt},
  txt/.style={font=\small,text=Ink,inner sep=0pt},
  note/.style={font=\footnotesize,text=Slate,inner sep=0pt},
  lab/.style={font=\small\bfseries,text=Ink,inner sep=0pt}
}
\begin{document}
\noindent
\begin{tikzpicture}[x=1mm,y=1mm,line cap=round,line join=round]
\path[use as bounding box] (0,0) rectangle (273,182);

% ---- Page heading (separate from all drawing annotations) ----
\node[anchor=west,font=\LARGE\bfseries,text=Ink,inner sep=0pt] at (0,175)
  {クリフの正面図と断面図};
\node[anchor=west,note] at (0,164)
  {本体前面に付く把持用の突起と，その下の有限な本体面を示す。};
\draw[Line,line width=.7pt] (0,157) -- (273,157);

% ---- Shared physical scale for the specified cross-sectional dimensions ----
% 1 physical centimetre = 2.3 drawing millimetres.
\def\sc{2.3}
\pgfmathsetmacro{\depth}{3*\sc}
\pgfmathsetmacro{\gripH}{5*\sc}
\pgfmathsetmacro{\underH}{15*\sc}
\def\gripTop{104}
\pgfmathsetmacro{\gripBottom}{\gripTop-\gripH}
\pgfmathsetmacro{\bodyBottom}{\gripBottom-\underH}
\def\bodyTop{130}

% ================================================================
% (a) Front view: observer on the player's side, looking at the face.
% ================================================================
\node[anchor=west,lab] at (0,149) {(a) 正面図};
\node[anchor=west,note] at (0,141) {選手側から本体前面を見る};

% Schematic supports (not cables and not an assumed pendulum mechanism).
\draw[Ink,line width=.9pt] (30,130) -- (30,136);
\draw[Ink,line width=.9pt] (88,130) -- (88,136);
\draw[Ink,line width=.5pt] (28.5,134.7) -- (31.5,135.5);
\draw[Ink,line width=.5pt] (86.5,134.7) -- (89.5,135.5);
\node[note] at (59,134.5) {上部から支持};

% Body and continuous front ledge.
\filldraw[edge,fill=Body] (13,\bodyBottom) rectangle (105,\bodyTop);
\filldraw[draw=Ledge,line width=.75pt,fill=LedgeFill]
  (13,\gripBottom) rectangle (105,\gripTop);
\draw[Ledge,line width=1.35pt] (13,\gripTop) -- (105,\gripTop);
\node[txt] at (50,119) {本体前面};
\node[font=\small\bfseries,text=Ledge,inner sep=0pt] at (50,98.2)
  {把持用の突起};
\node[txt] at (50,76) {突起の下に続く本体面};

% Front-view height dimensions (same scale as cross section).
\foreach \Y in {\gripTop,\gripBottom,\bodyBottom}{
  \draw[ext] (107,\Y) -- (116,\Y);
}
\draw[dim] (113,\gripBottom) -- (113,\gripTop);
\node[anchor=west,txt] at (117,98.25) {$5\,\mathrm{cm}$};
\draw[dim] (113,\bodyBottom) -- (113,\gripBottom);
\node[anchor=west,txt] at (117,75.25) {$15\,\mathrm{cm}$};

% Section plane A--A. The arrows point towards the right-hand sectional view.
\draw[Ink,dash pattern=on 4pt off 2pt on .7pt off 2pt,line width=.45pt]
  (77,133) -- (77,53);
\draw[-{Latex[length=1.8mm]},Ink,line width=.55pt] (73.5,133) -- (81,133);
\draw[-{Latex[length=1.8mm]},Ink,line width=.55pt] (73.5,53) -- (81,53);
\node[anchor=east,note] at (72.3,133) {A};
\node[anchor=east,note] at (72.3,53) {A};

% Longitudinal length: no invented numerical value.
\draw[ext] (13,56) -- (13,43);
\draw[ext] (105,56) -- (105,43);
\draw[dim] (13,45) -- (105,45);
\node[txt,fill=white,inner sep=1mm] at (59,45) {$L$：十分長い};
\node[anchor=north,align=center,note] at (59,35)
  {突起は長手方向に連続する。\\下端より下は開放されている。};

% ================================================================
% (b) Cross section: back (left), player/free space (right).
% The small ledge projects to the RIGHT of the body front face.
% ================================================================
\node[anchor=west,lab] at (139,149) {(b) A--A断面図};
\node[anchor=west,note] at (139,141) {長手方向 $L$ に直交する断面};

\def\back{166}
\def\front{184}
\pgfmathsetmacro{\tip}{\front+\depth}

% Solid section envelope. Internal construction is deliberately omitted.
\filldraw[edge,fill=Body] (\back,\bodyBottom) rectangle (\front,\bodyTop);
\begin{scope}
  \clip (\back,\bodyBottom) rectangle (\front,\bodyTop);
  \foreach \k in {-80,-74,...,80}{
    \draw[Line,line width=.25pt]
      (160,{\bodyBottom+\k}) -- (192,{\bodyBottom+\k+32});
  }
\end{scope}
\draw[edge] (\back,\bodyBottom) rectangle (\front,\bodyTop);
% Ledge section; the tip upper edge carries a 5 mm fillet (R5) so that it does not dig into the fingers.
\pgfmathsetmacro{\rfil}{0.5*\sc}
\filldraw[draw=Ledge,line width=.75pt,fill=LedgeFill]
  (\front,\gripBottom) -- (\tip,\gripBottom)
  {[rounded corners=\rfil mm] -- (\tip,\gripTop) -- (\tip-2,\gripTop)} -- (\front,\gripTop) -- cycle;
\draw[Ledge,line width=1.35pt] (\front,\gripTop) -- (\tip-\rfil,\gripTop);
\node[txt,fill=Body,inner sep=1mm] at (175,112.6) {本体};

% Support is diagrammatic, not a specified suspension mechanism.
\draw[Ink,line width=.9pt] (175,130) -- (175,136);
\draw[Ink,line width=.5pt] (173.5,134.7) -- (176.5,135.5);
\node[anchor=west,note] at (184,134.8) {上部から支持};

% Section height dimensions to the LEFT of the material; labels in free space.
\foreach \Y in {\gripTop,\gripBottom,\bodyBottom}{
  \draw[ext] (152,\Y) -- (164,\Y);
}
\draw[dim] (155,\gripBottom) -- (155,\gripTop);
\node[anchor=east,txt] at (151,98.25) {$5\,\mathrm{cm}$};
\draw[dim] (155,\bodyBottom) -- (155,\gripBottom);
\node[anchor=east,txt] at (151,75.25) {$15\,\mathrm{cm}$};

% 3 cm is measured from the FRONT FACE to the TIP only.
\draw[ext] (\front,105.8) -- (\front,125.5);
\draw[ext] (\tip,105.8) -- (\tip,125.5);
\draw[dim] (\front,122) -- (\tip,122);
\node[txt,fill=white,inner sep=.65mm] at (188,127) {$3\,\mathrm{cm}$};

% Hand profile: schematic centreline, entirely OUTSIDE the body and ledge.
% The finger pad lies over the ledge's upper surface. Neither hand nor wrist
% is drawn through the material. It is not an anatomical/kinematic model.
\draw[Hand,line width=1.6pt]
  (185.8,104.55) -- (188.6,104.55)
  .. controls (192.4,104.55) and (194.0,101.5) .. (194.0,98.5)
  -- (194.0,94.8)
  .. controls (194.0,91.0) and (188.2,90.5) .. (187.0,86.2)
  -- (186.2,81.7);
\filldraw[draw=Hand,line width=.8pt,fill=HandFill] (186.2,80.5) circle[radius=1.2mm];
\draw[Hand,line width=1.6pt] (186.2,79.3) -- (192.5,45.5);

% Contact surface. Small red segment is the actual front face, never the tip.
\draw[Contact,line width=1.2pt] (\front,73.5) -- (\front,86);

% Callouts occupy their own strip on the right; no paragraph over the figure.
\fill[Ledge] (184.7,104) circle[radius=.55mm];
\draw[leader] (184.7,104) -- (201,113) -- (212,113);
\node[anchor=west,lab,text=Ledge] at (215,114.4) {指を掛ける上面};
\node[anchor=west,note] at (215,108.3) {前面から $3\,\mathrm{cm}$ 張り出す。先端上縁は R$5\,\mathrm{mm}$};

\fill[Contact] (\front,76) circle[radius=.55mm];
\draw[leader] (\front,76) -- (207,85) -- (212,85);
\node[anchor=west,lab,text=Contact] at (215,88.4) {手首に干渉する面};
\node[anchor=north west,align=left,note] at (215,83.8)
  {体を奥へ振る際の，\\手首の移動を制限する。};

\draw[leader] (189.5,61.7) -- (207,65) -- (212,65);
\node[anchor=west,note,text=Hand] at (215,65) {手・前腕（模式）};

% Free space below the finite body. The arc is a direction indicator only.
\draw[<->,>=Latex,Ink,line width=.8pt]
  (165,42) .. controls (169,24) and (200,24) .. (205,42);
\node[anchor=west,lab] at (215,40.5) {下側は空間};
\node[anchor=north west,align=left,note] at (215,35.8)
  {本体下端より下で\\体を前後に振れる。};
\draw[leader] (195,28.7) -- (209,31);

% ---- Unambiguous modelling convention, separated from the views ----
\draw[Line,line width=.7pt] (0,18) -- (273,18);
\node[anchor=west,note,text=Ink] at (0,11.3)
  {寸法の解釈：$5\,\mathrm{cm}$ は突起の鉛直高さ，$15\,\mathrm{cm}$ は突起下面から本体下端まで。先端上縁の丸みは R$5\,\mathrm{mm}$（シミュレーションの \texttt{edge\_radius}）。};
\node[anchor=west,note] at (0,4.8)
  {形状参照：添付資料「クリフディメンションの背面飛び移りの物理学」1ページ。支持機構・本体奥行き・上部高さ・手の形は模式。};
\end{tikzpicture}
\end{document}
